\section{Why two scans are required}\label{sec:twoscans}

A natural idea for a faster protocol is to acquire only the full-resolution scan, fit a smooth
(e.g.\ low-order polynomial) $\Bone$ field to the whole image, and estimate $\Tone$, $\Mzero$ and the other
parameters voxel by voxel. The field has only a handful of coefficients, while the image has
$10^5$--$10^6$ voxels, so it seems that the data should over-determine it. This section shows that
they do not: the single-scan degeneracy of \cref{cor:onescan} is exact and local, and no
spatial model of $\Bone$ alone can remove it. A second acquisition in which the flip angle enters
differently is required; it can, however, be very short.

\paragraph{A toy model.} Suppose each voxel $v$ has an unknown amplitude $a_v$ and the scanner has an
unknown gain $b$ that varies smoothly (here, a constant), and one acquisition measures
$y_v=a_v\,b$. However many voxels there are, $b$ cannot be found: replacing $(a_v,b)$ with
$(a_v/\kappa,\kappa b)$ leaves every $y_v$ unchanged. The smoothness of $b$ does not help, because each
$a_v$ is free and absorbs the change locally. A second acquisition that depends on $b$
differently, for example $y'_v=a_v\,b^2$, breaks the degeneracy in every voxel: $b=y'_v/y_v$. The
single MPME scan is this situation with $a_v\leftrightarrow(\Mzero,\Tone)$ and $b\leftrightarrow\Bone$. The data
determine only the flip angle relative to the Ernst angle, $\xi=\tan(\Bone\alpha/2)/\tan(\alpha_E/2)$, and the
amplitude $\Mzero\tan(\alpha_E/2)$ (\cref{prop:ernst,tab:onescan}). For every trial $B_1^{+\prime}$ the voxel's own
$\Tone'$ moves its Ernst angle to restore $\xi$, and its own $\Mzero'$ restores the amplitude
(\eqref{eq:family}, \cref{tab:family}). Sharing $\Bone$ with other voxels does not change this, because
each voxel compensates with its own parameters (\cref{sec:shared}).

\paragraph{Counting.} With $N$ voxels and a $\Bone$ model with $p$ coefficients, one scan supplies
$2N$ independent equations for these unknowns (the two invariants per voxel), while the unknowns
number $2N+p$ ($\Mzero$ and $\Tone$ in each voxel, plus the field). The system is underdetermined by
exactly $p$ for every $N$, so adding voxels never closes the gap. A second flip angle adds $N$
equations ($\xi_2$ in each voxel, which with $\xi_1$ gives \eqref{eq:secondscan}), giving $3N\ge2N+p$. With a smooth $\Bone$
model only about $p$ of the new equations are needed to fix the field, which is why the second
scan can be short (\cref{sec:design}).

\subsection{Any $\Bone$ field fits a single scan exactly}

\begin{proposition}[Whole-image degeneracy of a single scan]\label{prop:singlescan}
Consider one scan with flip angle $\alpha^{\rm nom}$ and data generated by
$(\Mzero(\mathbf r),\Bone(\mathbf r),\Tone(\mathbf r))$ and arbitrary $\Ttwo$, $\Ttwos$, $\dw$, $\phi_0$. Let
$B_1^{+\prime}(\mathbf r)$ be \emph{any} field satisfying
\begin{equation}
  \cos\bigl(B_1^{+\prime}(\mathbf r)\,\alpha^{\rm nom}\bigr)>-u(\mathbf r)\qquad\text{for every voxel},
  \label{eq:admissible}
\end{equation}
where $u(\mathbf r)$ is the invariant \eqref{eq:uv} of the true parameters. Then there are fields
$\Tone'(\mathbf r)$ and $\Mzero'(\mathbf r)>0$ such that $(\Mzero',B_1^{+\prime},\Tone')$, with the same $\Ttwo$, $\Ttwos$, $\dw$ and
$\phi_0$ (up to $\phi_0\to\phi_0+\pi$), produce exactly the same complex signal in every voxel, every
pathway and every echo.
\end{proposition}
\begin{proof}
Fix a voxel and write $c'=\cos(B_1^{+\prime}\alpha^{\rm nom})$. Inverting the M\"obius map \eqref{eq:uv}
in $E_1$, the value
\[
  E_1'=\frac{u+c'}{1+uc'}
\]
gives $u(B_1^{+\prime}\alpha^{\rm nom},E_1')=u$. Since $|u|<1$ and $|c'|<1$, the denominator is positive;
$E_1'<1$ is equivalent to $(1-u)(1-c')>0$, which always holds; and $E_1'>0$ is equivalent to
\eqref{eq:admissible}. Hence $\Tone'=-\TR/\log E_1'$ is a valid relaxation time. By \cref{thm:uv} the
signal is $\Mzero\,v\,G_k(u,E_2,\dw)$ times echo factors that do not involve $\Mzero$, $\Bone$ or $\Tone$; with
$u$ unchanged, choosing $\Mzero'=\Mzero\,|v|/|v'|$ makes $\Mzero'|v'|=\Mzero|v|$, and if $\sgn v'\ne\sgn v$ the
sign is absorbed by $\phi_0\to\phi_0+\pi$ (a single scan has a single phase). The voxel's signal is
therefore reproduced exactly. The construction is independent from voxel to voxel, so it
applies to any field $B_1^{+\prime}(\mathbf r)$.
\end{proof}

Condition \eqref{eq:admissible} is no restriction in practice. For the published scan 1
($15^\circ$, $\TR=25$\,ms), $u=0.08$ (WM), $0.30$ (GM) and $0.69$ (CSF), so any field with
$B_1^{+\prime}<6.3$, $7.2$ and $8.9$, respectively, is admissible---every physically plausible field.
The statement holds at every flip angle because \cref{thm:uv} is exact; it is not a
small-angle approximation. In the small-angle limit the compensation reduces to the
scaling of \cref{prop:scaling}, $\Tone'\propto(B_1^{+\prime})^{-2}$ and $\Mzero'\propto(B_1^{+\prime})^{-1}$.

\begin{corollary}[A parametric $\Bone$ model has no information from one scan]\label{cor:parametric}
Let $\Bone(\mathbf r)=\sum_{j=1}^p\beta_j\phi_j(\mathbf r)$ for any basis $\{\phi_j\}$ (polynomials, splines,
\dots), with $\Tone$, $\Mzero$ (and $\Ttwo$, $\Ttwos$, $\dw$, $\phi_0$) free in every voxel. From one scan, the
Fisher information about $\beta$, after the per-voxel parameters are accounted for (profiled),
is the zero matrix, and the CRLB on every $\beta_j$ is infinite.
\end{corollary}
\begin{proof}
At each voxel $v$, \cref{cor:onescan} shows that the three Jacobian columns of $S_v$ with
respect to $\log\Mzero$, $\log\Bone$ and $\log\Tone$ span at most two dimensions; the columns for $\log\Mzero$
and $\log\Tone$ are generically independent, so
\[
  \frac{\partial S_v}{\partial\log\Bone}=a_v\,\frac{\partial S_v}{\partial\log\Tone}+b_v\,\frac{\partial S_v}{\partial\log\Mzero}
\]
for some scalars $a_v,b_v$. The Jacobian column of a coefficient $\beta_j$ is the stack over voxels of
$\phi_j(\mathbf r_v)\,\Bone(\mathbf r_v)^{-1}\,\partial S_v/\partial\log\Bone$, which therefore lies in the span of the per-voxel
$\log\Tone$ and $\log\Mzero$ columns. The profiled information of $\beta$ is
$\Fi_{\beta\beta}-\Fi_{\beta\eta}\Fi_{\eta\eta}^{-1}\Fi_{\eta\beta}$ with $\eta$ the per-voxel parameters; it is the Gram matrix
of the components of the $\beta$-columns orthogonal to the $\eta$-columns, which are zero.
\end{proof}

The number of voxels does not matter: each voxel contributes exactly zero information about
$\Bone$, and a sum of zeros is zero. Smoothness reduces the number of $\Bone$ unknowns, but not the
per-voxel freedom of $\Tone$ and $\Mzero$ to compensate.

\paragraph{Numerical confirmation.} \Cref{tab:singlescan} refits noise-free scan-1 data of the
synthetic phantom (4735 voxels across the slice) with three $\Bone$ fields held fixed: the true
field, the true field scaled by $1.10$, and the true field times a smooth tilt and curvature.
All three fit to machine precision; the $\Tone$ maps differ by $\mathrm d\log\Tone/\mathrm d\log\Bone=-2.03$ and
$-2.02$, i.e.\ a $10\%$ $\Bone$ error becomes a uniform $-19\%$ $\Tone$ error, and the smooth tilt becomes a
spurious $\Tone$ gradient from $-22\%$ to $+27\%$ across the slice. The profiled information of the six
coefficients of a degree-2 polynomial over 60 voxels with varied tissue is zero to rounding
(largest eigenvalue $4\times10^{-16}$) for one scan, and positive definite (smallest eigenvalue
$0.89$, largest $857$) for the two scans of the published protocol.

\begin{table}[!htbp]
\centering
\caption{Scan-1 data refitted with fixed $\Bone$ fields (noise-free, 4735 voxels).}
\label{tab:singlescan}
\small
\begin{tabular}{lccc}
\toprule
$\Bone$ field used & max relative residual & $\log(\hat\Tone/\Tone)$: median [range] & $\mathrm d\log\Tone/\mathrm d\log\Bone$\\
\midrule
true & $4\times10^{-16}$ & $0.000$ & --\\
true $\times1.10$ & $6\times10^{-15}$ & $-0.193$ [$-0.194$, $-0.192$] & $-2.03$\\
true $\times(1+0.15x-0.10y^2)$ & $8\times10^{-15}$ & $+0.055$ [$-0.225$, $+0.274$] & $-2.02$\\
\bottomrule
\end{tabular}
\end{table}

\subsection{Voxels that share $\Bone$}\label{sec:shared}

The extreme case of a smooth field is a constant one: a group of voxels with the same $\Bone$ and
otherwise different tissue. Joint estimation with one shared $\Bone$ still leaves an exact
one-parameter family, but within the group it preserves the ratios of $\Mzero$ exactly and the
ratios of $\Tone$ almost exactly.

\begin{proposition}[Voxels sharing $\Bone$]\label{prop:shared}
Let voxels $j=1,\dots,N$ share $\Bone$, with arbitrary $M_{0,j}$, $T_{1,j}$, $T_{2,j}$, $T_{2,j}^{*}$, $\Delta\omega_j$, $\phi_{0,j}$,
and let $\xi_j$, $\zeta_j=\tan(\alpha_{E,j}/2)$ be their scan-1 quantities (\cref{prop:ernst}). For every
$B_1^{+\prime}$ with $|\tan(B_1^{+\prime}\alpha^{\rm nom}/2)|<\min_j|\xi_j|$, the parameters
\begin{equation}
  \zeta_j'=\frac{\tan(B_1^{+\prime}\alpha^{\rm nom}/2)}{\xi_j},\qquad
  T_{1,j}'=\frac{\TR}{2\operatorname{artanh}\zeta_j'^2},\qquad
  M_{0,j}'=M_{0,j}\frac{\zeta_j}{\zeta_j'},
  \label{eq:sharedfamily}
\end{equation}
with all other parameters unchanged, reproduce the scan-1 signal of every voxel exactly. Along
this family, for all $j,l$,
\begin{equation}
  \frac{M_{0,j}'}{M_{0,l}'}=\frac{M_{0,j}}{M_{0,l}},\qquad
  \frac{\tanh(\TR/2T_{1,j}')}{\tanh(\TR/2T_{1,l}')}=\frac{\tanh(\TR/2T_{1,j})}{\tanh(\TR/2T_{1,l})}
  \label{eq:sharedratios}
\end{equation}
exactly, and $T_{1,l}'/T_{1,j}'=(T_{1,l}/T_{1,j})\bigl(1+O\bigl((\TR/2T_{1,\min}')^2\bigr)\bigr)$.
\end{proposition}
\begin{proof}
For each voxel, \eqref{eq:sharedfamily} is \eqref{eq:family} with the common $B_1^{+\prime}$; it keeps $\xi_j$
and $M_{0,j}\zeta_j$, hence the voxel's signal (\cref{prop:singlescan}), and the condition on $B_1^{+\prime}$ makes
every $\zeta_j'<1$. The voxels' likelihood terms share no other parameter, so the joint data are
reproduced. Dividing, $\zeta_j'/\zeta_l'=\xi_l/\xi_j$ does not depend on $B_1^{+\prime}$, and it equals $\zeta_j/\zeta_l$ at
$B_1^{+\prime}=\Bone$; since $\zeta^2=\tanh(\TR/2\Tone)$ this is the second identity. Then
$M_{0,j}'/M_{0,l}'=(M_{0,j}\zeta_j/M_{0,l}\zeta_l)(\zeta_l'/\zeta_j')=M_{0,j}/M_{0,l}$. Finally
$\operatorname{artanh}y=y\,(1+y^2/3+O(y^4))$ with $y=\zeta^2\approx\TR/2\Tone$ gives the $\Tone$ ratio.
\end{proof}

In words: changing the shared $\Bone$ multiplies $\tan(\Bone\alpha/2)$ by a factor $\kappa$, and every voxel keeps
its own $\xi_j$ by multiplying its own $\tan(\alpha_{E,j}/2)$ by the same $\kappa$. All Ernst angles therefore
slide together, all $\Tone$ values change by about $\kappa^{-2}$ and all $\Mzero$ values by exactly
$\kappa^{-1}$. The shared parameter is never tested, because no voxel's data
involve another voxel's parameters. This is \cref{cor:parametric} with a single constant basis
function: $2N$ equations for $2N+1$ unknowns.

\Cref{tab:sharedfamily} shows the family for a WM voxel ($\Mzero=0.69$, $\Tone=850$\,ms) and a GM voxel
($\Mzero=0.82$, $\Tone=1350$\,ms) that share $\Bone=1$ and differ in every other parameter as well. Both
voxels reproduce their scan-1 signals to machine precision for every $B_1^{+\prime}$ from $0.7$ to $1.5$, while
the absolute $\Tone$ values change by a factor of $4.7$. The ratio $\Mzero^{\rm GM}/\Mzero^{\rm WM}$ stays at
$0.82/0.69=1.188406$ to all printed digits, and $\Tone^{\rm GM}/\Tone^{\rm WM}$ stays between $1.5882$ and
$1.5885$ (true value $1.5882$). Scan 2 again rejects every member but the truth.

\begin{table}[!htbp]
\centering
\caption{Two voxels with the same $\Bone$ (truth $1$) fitted jointly from scan 1: the shared-$\Bone$
family \eqref{eq:sharedfamily}. Misfit columns as in \cref{tab:family}, maximum over the two
voxels. Generated by \texttt{scripts/single\_scan\_family.py}.}
\label{tab:sharedfamily}
\small\setlength{\tabcolsep}{4.5pt}
\begin{tabular}{ccccccccc}
\toprule
& \multicolumn{3}{c}{$\Tone'$ (ms)} & \multicolumn{3}{c}{$\Mzero'$} & \multicolumn{2}{c}{misfit}\\
\cmidrule(lr){2-4}\cmidrule(lr){5-7}\cmidrule(lr){8-9}
$B_1^{+\prime}$ & WM & GM & GM/WM & WM & GM & GM/WM & scan 1 & scan 2\\
\midrule
0.7 & 1745 & 2771 & 1.5882 & 0.989 & 1.175 & 1.188406 & $1\times10^{-15}$ & 0.86\\
0.8 & 1334 & 2118 & 1.5882 & 0.864 & 1.027 & 1.188406 & $2\times10^{-15}$ & 0.70\\
0.9 & 1052 & 1670 & 1.5882 & 0.768 & 0.912 & 1.188406 & $3\times10^{-15}$ & 0.43\\
1.0 & 850 & 1350 & 1.5882 & 0.690 & 0.820 & 1.188406 & $2\times10^{-16}$ & 0.00\\
1.1 & 701 & 1113 & 1.5883 & 0.627 & 0.745 & 1.188406 & $5\times10^{-16}$ & 1.42\\
1.2 & 587 & 933 & 1.5883 & 0.574 & 0.682 & 1.188406 & $4\times10^{-16}$ & 2.00\\
1.5 & 372 & 591 & 1.5885 & 0.457 & 0.543 & 1.188406 & $1\times10^{-15}$ & 1.27\\

\bottomrule
\end{tabular}
\end{table}

\paragraph{A joint fit can look as if it worked.} \Cref{tab:sharedfit} fits the WM and GM voxels of
\cref{tab:sharedfamily} jointly, with one shared $\Bone$, to noisy data ($\sigma=10^{-3}$, one
realisation). The method is a block Levenberg--Marquardt with exact Jacobians, started from a
generic point ($\Mzero=0.8$, $\Tone=1000$\,ms, $\Ttwo=80$\,ms, $\Ttwos=56$\,ms) at five values of $\Bone$. From
scan 1 alone, every fit reaches the same cost, $22.9687\,\sigma^2$, but the final $\Bone$ depends on the
start (from $0.90$ to $1.00$): the optimiser stops at the first point of the solution curve it
reaches. A shared-$\Bone$ fit that starts near the truth and ends near it is therefore no
evidence that the data determine $\Bone$. The other diagnostics agree:
\begin{itemize}
\item \emph{Profile.} Over $B_1^{+\prime}\in[0.6,1.6]$ the profile cost is constant to $5\times10^{-14}\sigma^2$, and to
  $5\times10^{-13}\sigma^2$ for a set of 20 voxels with random tissue parameters.
\item \emph{Fisher information.} The joint Fisher matrix has one zero eigenvalue (below
  $4\times10^{-16}$ after normalisation), whose eigenvector is the family tangent
  \eqref{eq:familyslopes} to $10^{-13}$. The profiled information on $\log\Bone$ is zero to rounding.
\end{itemize}
With both scans, every start up to $1.1$ reaches the minimum at $\Bone=1.005$ (20 voxels: $1.0004$).
The curvature of the profile there gives $\mathrm{SD}(\log\Bone)=2.14\times10^{-3}$, equal to the
information at the estimate (20 voxels: $5.52\times10^{-4}$ against the bound $5.54\times10^{-4}$). Starts at
$1.2$ and above end on the bound $540/330$ at a much higher cost, which a multi-start fit discards.
Without the bound they reach the complex alias $2.016$ at exactly the minimal cost. Because
$f$ in \eqref{eq:secondscan} is the same in every voxel, sharing $\Bone$ does not remove the alias.

\begin{table}[!htbp]
\centering
\caption{Joint fit of the WM and GM voxels with a shared $\Bone$ (truth $1$), $\sigma=10^{-3}$, from a
generic start at five values of $\Bone$; fitting box $\Bone\in[\E^{-1},540/330]$. Generated by
\texttt{scripts/shared\_b1\_check.py}.}
\label{tab:sharedfit}
\small
\begin{tabular}{ccccc}
\toprule
& \multicolumn{2}{c}{scan 1 only} & \multicolumn{2}{c}{both scans}\\
\cmidrule(lr){2-3}\cmidrule(lr){4-5}
start $\Bone$ & final $\Bone$ & cost$/\sigma^2$ & final $\Bone$ & cost$/\sigma^2$\\
\midrule
0.6 & 0.8966 & 22.9687 & 1.0050 & 76.0635\\
0.8 & 0.9538 & 22.9687 & 1.0050 & 76.0635\\
1.0 & 0.9821 & 22.9687 & 1.0050 & 76.0635\\
1.3 & 0.9981 & 22.9687 & 1.6364 & $3.9\times10^{4}$\\
1.6 & 1.0032 & 22.9687 & 1.6364 & $3.9\times10^{4}$\\

\bottomrule
\end{tabular}
\end{table}

So one scan with a shared (or smooth) $\Bone$ determines \emph{relative} $\Tone$ and $\Mzero$ maps
within a region of constant $\Bone$, to within $O((\TR/2\Tone)^2)$, but not their common scale: an
error $\delta\log\Bone$ multiplies every $\Tone$ by about $\E^{-2\delta\log\Bone}$ and every $\Mzero$ by about
$\E^{-\delta\log\Bone}$. The scale can be fixed only by information from outside the scan, for example a
reference voxel whose $\Tone$ is assumed known: then $\Bone=(2/\alpha^{\rm nom})\arctan(\xi_{\rm ref}\zeta_{\rm ref})$.
A relative error $\delta$ in the assumed reference $\Tone$ then becomes a $\Bone$ error of about $-\delta/2$
and an error of about $\delta$ in every $\Tone$ of the region. This is a prior of the kind discussed
next, and the accuracy of the maps is that of the assumption.

\subsection{What a prior would do instead}

Identifiability can be restored by constraining $\Tone$ or $\Mzero$ spatially, for example $\Tone$
constant within each tissue class, or $\Mzero=C\cdot\mathrm{PD}$ with a smooth $C$ and class-constant $\mathrm{PD}$. The
field is then determined by the prior, not by the data: any departure of the true $\Tone$ from the
assumed uniformity is absorbed into $\Bone$. By the compensation law $\delta\log\Tone=-2\,\delta\log\Bone$, a real
regional $\Tone$ difference $\delta$ within a tissue is converted into a $\Bone$ change of about $-\delta/2$ and
removed from the $\Tone$ map. For quantitative mapping, whose purpose is to detect such
differences, this is circular. Priors are useful on top of a second acquisition (e.g.\ the
hierarchical model of \cref{sec:future}), not in place of one.

\subsection{Physics outside the ideal model}\label{sec:beyond}

The degeneracy belongs to the ideal model: instantaneous pulses, one flip angle per voxel and
no diffusion. Real data contain effects that make the signal depend on the flip angle in other
ways, and each could in principle give a single scan some information about $\Bone$.
\Cref{tab:beyond} adds each effect to the model and gives the scan-1 bound on $\log\Bone$, with every
other parameter free, against $2.2$ ($\sigma/\Mzero$ units, WM) for the ideal two-scan protocol.

\begin{table}[!htbp]
\centering
\caption{Information about $\Bone$ in scan 1 ($15^\circ$) when the ideal model is extended: CRLB on
$\log\Bone$ in units of $\sigma/\Mzero$, all other parameters free ($\log\Mzero$, $\log\Tone$, $\log\Ttwo$, $\log\Ttwos$, $\dw$,
$\phi_0$, plus the extra parameter where one is unknown). ``Exactly degenerate'': zero
information to rounding. Generated by \texttt{scripts/beyond\_ideal\_model.py}.}
\label{tab:beyond}
\small\setlength{\tabcolsep}{4pt}
\begin{tabular}{p{0.30\linewidth}p{0.24\linewidth}cc}
\toprule
Effect & Assumption & WM & GM\\
\midrule
Ideal model & reference: both scans & 2.21 & 2.75\\
 & scan 1 & exactly degenerate & exactly degenerate\\
\midrule
Gaussian slice profile, $|z|\le3\sigma$ & shape known & $8.4\times10^{5}$ & $5.9\times10^{5}$\\
 & width unknown & exactly degenerate & exactly degenerate\\
Hann-sinc slice profile, TBW 4 & shape known & $1.2\times10^{6}$ & $8.3\times10^{5}$\\
 & width unknown & exactly degenerate & exactly degenerate\\
Slab-edge ramp, $p\in[0.9,1]$ & shape known & $1.2\times10^{7}$ & $1.2\times10^{7}$\\
 & width unknown & exactly degenerate & exactly degenerate\\
\midrule
Relaxation during the pulse, on resonance & $\tau=1$\,ms & exactly degenerate & exactly degenerate\\
 & $\tau=2$\,ms & exactly degenerate & exactly degenerate\\
Off-resonance during the pulse, $\tau=1$\,ms & $\dw/2\pi=8$\,Hz & $5.0\times10^{4}$ & $5.3\times10^{4}$\\
 & $\dw/2\pi=48$\,Hz & $8.4\times10^{3}$ & $8.9\times10^{3}$\\
 & $\dw/2\pi=160$\,Hz & $2.7\times10^{3}$ & $2.8\times10^{3}$\\
\midrule
Diffusion, $D=0.7\times10^{-3}$\,mm$^2$/s & $q=2\pi/1$\,mm, $D$ known & $4.8\times10^{3}$ & $3.7\times10^{3}$\\
 & $q=2\pi/1$\,mm, $D$ free & exactly degenerate & exactly degenerate\\
 & $q=2\pi/0.5$\,mm, $D$ known & $1.5\times10^{3}$ & $1.2\times10^{3}$\\
 & $q=2\pi/0.5$\,mm, $D$ free & exactly degenerate & exactly degenerate\\

\bottomrule
\end{tabular}
\end{table}

\begin{itemize}
\item \emph{Flip-angle distribution within the voxel} (slice or slab profile), each position in
  its own steady state. With the profile shape known exactly, the bound is finite but
  $5.9\times10^5$--$1.2\times10^7$, five to seven orders of magnitude above the two-scan value. With
  the profile width unknown, the degeneracy is exact again.
\item \emph{Relaxation during an on-resonance pulse} ($\tau=1$--$2$\,ms). The signal changes by
  $0.15$--$0.8\%$, but the degeneracy remains exact. As for an ideal pulse, the pathway amplitudes
  keep a one-ratio structure: $a_k/a_0=r^k$ for $k\ge0$ and $a_{-1}/a_0=(r-E_2)/(1-E_2r)$ (verified to
  $10^{-16}$).
\item \emph{Off-resonance during a $1$\,ms pulse} (8--160\,Hz). Bound $2.7\times10^3$--$5.3\times10^4$, carried by
  small pathway-dependent changes that unmodelled gradient delays or eddy currents would
  easily mask.
\item \emph{Diffusion in the unbalanced gradient} ($D=0.7\times10^{-3}$\,mm$^2$/s, one cycle of dephasing per
  $1$ or $0.5$\,mm). With $D$ known the bound is $1.2$--$4.8\times10^3$; with $D$ free the degeneracy is
  exact.
\end{itemize}

A simple count explains the pattern. The echo decay and phases of a scan determine $\Ttwo$, $\Ttwos$,
$\dw$ and $\phi_0$, which leaves three real pathway amplitudes for $(\Mzero,\Bone,\Tone)$. In the ideal model
they depend on only two combinations. An effect that adds structure supplies a weak third,
unless it brings an unknown of its own (a profile width, $D$), which absorbs it exactly.

Even in the best cases this information is of no practical use. Pooled through a smooth
field \eqref{eq:smoothsd} with $p=35$ and $N=10^5$ at SNR 300, the bounds of \cref{tab:beyond} give
$\mathrm{SD}(\log\Bone)\approx0.09$ (diffusion, $0.5$\,mm, $D$ known) to $0.17$ (160\,Hz), against $1.4\times10^{-4}$ for
the two-scan protocol. Even these values require the extra physics to be modelled exactly,
and the profile shows how strongly a model error biases the answer. If the true Gaussian
profile is $2\%$ narrower than assumed, the scan-1 fit converges exactly (residual $10^{-16}$) to
$\Bone=1.47$, $\Tone=386$\,ms for WM (truth $1$ and $850$\,ms); $5\%$ narrower gives $\Bone=1.89$; for the ramp,
$2\%$ gives $\Bone=1.73$; for the Hann-sinc profile the bias has the opposite sign ($\Bone=0.93$). If the
true Gaussian or ramp profile is wider than assumed, no exact solution exists. The fit then
runs to $\Bone\to0$, $\Tone\to\infty$ with a residual of $10^{-7}$--$10^{-4}$ of the signal, far below the noise. A single scan cannot be rescued by modelling more physics, any more than by a
smoother $\Bone$.

\subsection{What a second acquisition does}

With a second scan of the same $\TR$ and a different nominal flip angle, $\alpha_2=r\,\alpha_1$ with a
common $\Bone$, a voxel provides $u_1$, $u_2$ and, by \cref{lem:absv}, $\Mzero\sqrt{\tanh(\TR/2\Tone)}$: three numbers
for the three unknowns $(\Mzero,\Bone,\Tone)$. The construction of \cref{prop:singlescan} now fails:
$E_1'$ must satisfy $u(B_1^{+\prime}\alpha_1,E_1')=u_1$ and $u(B_1^{+\prime}\alpha_2,E_1')=u_2$ simultaneously, which are the two
equations \eqref{eq:aliaseq} in two unknowns. In Ernst-angle form, the ratio $\xi_2/\xi_1$ of the two scans
is the function $f(\Bone)$ of \eqref{eq:secondscan}, in which $\Tone$ and $\Mzero$ cancel. Its solutions are isolated: the per-voxel null
curve collapses to discrete points, the true parameters and the aliases of
\cref{prop:equiv}, which the FID phase and the range bound remove. Locally the information
is nonsingular (\cref{tab:crlb}), and the profiled information of a parametric field is
positive definite.

The single-scan family is the familiar $\Bone$ sensitivity of variable-flip-angle (VFA, DESPOT1)
$\Tone$ mapping. With $\Bone$ known, one MPME scan determines $\Tone$ and $\Mzero$, as a complete VFA
experiment does; with $\Bone$ unknown, it inherits the VFA error law $\delta\ln\Tone\approx-2\,\delta\ln\Bone$,
$\delta\ln\Mzero\approx-\delta\ln\Bone$ (exactly \eqref{eq:familyslopes}). Its pathways and echoes do not change this.
Spoiled VFA data give one amplitude per angle: two angles cannot determine $\Bone$ (two numbers
for three unknowns), and more angles determine it only weakly, because on the axis of
\cref{fig:logtan} their samples move with $\Bone$ at nearly equal rates $\alpha_i/\sin\alpha_i$ and
translate almost together with the bump. Each MPME scan also reports its offset from the
bump, so two scans suffice in principle; but two small angles are again nearly degenerate
(\cref{fig:logtan}b), and the published design relies on a second angle whose sample moves at
a very different rate.

Two variants achieve the same in principle: a second scan with the same flip angle and a
different $\TR$ (then $E_1$ changes while $\alpha$ does not, the principle of actual-flip-angle imaging),
or transient (non-steady-state) data, to which \cref{thm:uv} does not apply. The published
design uses a second flip angle near $360^\circ$, whose $\Bone$ lever arm \eqref{eq:lever} is large.

\subsection{How much second scan is needed}\label{sec:design}

Once each voxel carries some information about $\Bone$, pooling it through a smooth field is
effective, and the second scan can be short and coarse. If scan 2 covers a fraction $r$ of
$k$-space and both scans are reduced to that resolution, each low-resolution voxel has about
$1/\sqrt r$ times the per-voxel SNR of a full-resolution voxel, and there are $rN$ of them for $N$
full-resolution voxels. A parametric field with $p$ coefficients fitted to them has
\begin{equation}
  \mathrm{SD}(\log\hat B_1^{+})\approx\mathrm{CRLB}_{\Bone}\cdot\frac{\sigma}{\Mzero}\sqrt r\cdot\sqrt{\frac{p}{rN}}
  =\mathrm{CRLB}_{\Bone}\cdot\frac{\sigma}{\Mzero}\sqrt{\frac pN},
  \label{eq:smoothsd}
\end{equation}
independent of $r$ to this order. With $\mathrm{CRLB}_{\Bone}\approx2.5$, $p=35$ and $N=10^5$--$10^6$ this is of
order $10^{-4}$ at SNR 300: noise in the smoothed field is negligible, and its accuracy is limited
by bias (truncation ringing, partial volume, approach to steady state, model mismatch), which
grows as the coverage shrinks. \Cref{tab:design} summarises the design implications.

\begin{table}[!htbp]
\centering
\caption{Design of the second acquisition (bounds per unit $\sigma/\Mzero$, white-matter-like tissue,
$\Bone$ over $[0.8,1.2]$; median and worst case over $\Bone$).}
\label{tab:design}
\small
\begin{tabular}{p{0.22\linewidth}p{0.30\linewidth}p{0.38\linewidth}}
\toprule
Choice & Recommendation & Evidence\\
\midrule
$\alpha_2^{\rm nom}$ & $\approx330^\circ$ & $\Bone$ bound median $2.8$ at $330^\circ$ vs $3.5$ ($320^\circ$), $3.1$ ($340^\circ$), $7.0$ ($290^\circ$) with $\alpha_1=15^\circ$; the worst case ($11$--$13$) hardly depends on $\alpha_2$, so a $360^\circ$ crossing inside the head is unavoidable\\
$\alpha_1^{\rm nom}$ & $\approx20^\circ$ & improves both stages: $\Bone$ bound median $2.44$, worst $8.5$ (vs $2.77$, $11.7$ at $15^\circ$); scan-1 $\Tone$ bound with $\Bone$ known WM $46$, GM $52$ (vs $53$, $54$)\\
Scan-2 readout & same three windows, same $\TR$ & a one-window scan 2 with $\TR=12$\,ms is less efficient at equal duration (median $3.8$ vs $2.8$): fewer echoes constrain $\Ttwo$, $\Ttwos$ less\\
Scan-2 coverage & small, e.g.\ $24\times24$ central $k_y$--$k_z$ lines ($\approx8$\,mm) & \eqref{eq:smoothsd}: coverage barely affects the smoothed noise; bias limits it (apodise both scans identically; robust, uncertainty-weighted fit)\\
Preparation & $\ge5\Tone/\TR$ dummy repetitions & a short scan 2 must start in steady state\\
Phase & no re-shim or receiver change between scans & complex ML and the FID-sign test use the relative phase\\
\bottomrule
\end{tabular}
\end{table}

With $24\times24$ lines the second scan takes about 14\,s plus 7\,s of preparation, under $3\%$ of a
whole-brain exam, instead of about 48\,s in the published design. The claim that coverage can be
reduced this far without bias is argued from \eqref{eq:smoothsd} and the phantom results of
\cref{sec:phantom}, not yet demonstrated by a coverage sweep. The two-stage reconstruction
that uses this scan is given in \cref{alg:twostage} and evaluated in \cref{sec:phantom}.
